%!TEX root = ../../thesis.tex

\section{Relation to Prior Work}\label{sec:mdp-prior-work}

% TODO: Position both formulations against the survey in \cref{chap:related-work}.
% This section is common to \cref{sec:input-actions,sec:corpus-actions} because the
% two formulations descend from the same literature and the comparison table below
% would otherwise be built twice.
%
% Open with what BOTH inherit, stated once: the MDP formulation of fuzzing itself ---
% states as input fragments, actions as rewrite rules or symbols, reward from runtime
% coverage \cite{bottinger_deep_2018}. This thesis extends the canonical formulation
% rather than proposing a new one; say that plainly.
%
% Then what both depart from, also once: the learner. Model-based planning over a
% learned latent model, rather than a model-free learner tuning one knob of a
% conventional fuzzer. Every difference in the two subsections below is downstream of
% that single choice --- which is the argument the chapter has been building since
% \cref{sec:mdp-constraints}.

%!TEX root = ../../thesis.tex

\subsection{The Input-Level Formulation}\label{sec:input-prior-work}

% TODO: Points to make (the shared inheritance and the shared departure in the
% learner are stated above --- do not repeat them here):
%   - closest in spirit to the protocol/state-selection line: the agent drives the
%     session rather than advising a random mutator \cite{borcherding_bandits_2023},
%     and to the stateful greybox fuzzers those papers baseline against (AFLNet,
%     SGFuzz),
%   - differs from \cite{bottinger_deep_2018} in what an action *is*: a rewrite rule
%     applied to an existing seed there, an input symbol here,
%   - the honest limitation: generating an input symbol by symbol is not what a
%     production fuzzer does, which is precisely the motivation for
%     \cref{sec:corpus-actions}.

%!TEX root = ../../thesis.tex

\subsection{The Corpus-Level Formulation}\label{sec:corpus-prior-work}

% TODO: This formulation is a direct descendant of \cite{bottinger_deep_2018}, so give
% it a proper comparison table --- one row per MDP component, one column per approach.
% Rows: RL algorithm, state, action, mutation randomness, reward, corpus, horizon,
% target and instrumentation, reported result.
%
% What differs, beyond the shared departure in the learner stated above, and why each
% difference is forced rather than cosmetic:
%   1. model-based planning instead of model-free DQN, which forces the determinism
%      of \cref{sec:corpus-determinism},
%   2. a learned corpus --- selection *and* retention/eviction --- instead of a single
%      seed; note that the extension to a multi-seed corpus is mentioned there as future work with
%      a *random* choice of seed, never learned and never evaluated. This is the novel
%      core of this formulation,
%   3. a long horizon with seed lineage instead of a reset after every mutation.
%
% Also inherit the cautions, and state them before presenting any result:
%   - the reported gain over a random baseline was modest and called preliminary, so
%     calibrate expectations about beating random fuzzing,
%   - learned control of the mutation *offset* did not beat a random offset there,
%     which puts the learned position factor of \cref{sec:corpus-action-space} under
%     suspicion --- \cref{sec:results-ablations} should test whether it earns its
%     contribution to the branching factor.
%
% Finally, relate to the seed-scheduling line of \cref{sec:rlfuzz-prior-work}
% (bandit and MCTS-over-seed-tree approaches): those learn scheduling around a
% conventional fuzzer, whereas here selection, mutation and retention are one joint
% action in a single MDP.

